% Declaración de uso de IA: contenido sincronizado con la fuente Markdown.
\usoIA{Apertura}{Asistente LLM y Codex}{Integración del capítulo de aseguramiento de calidad}{Alto}{Ninguno}{Carlos Abarza: lectura del alcance, aclaración de términos y revisión de consistencia de los criterios de calidad declarados.}
\usoIA{9.1 Plan de Calidad}{Asistente LLM y Codex}{Conciliación de cobertura contractual y de ramas, métricas y gates}{Alto}{Ninguno}{Carlos Abarza: revisión de métricas, umbrales, fuentes y distinción entre mínimos contractuales y compromisos adicionales.}
\usoIA{9.2 Estrategia de Aseguramiento de Calidad}{Asistente LLM y Codex}{Desarrollo de estrategia y controles; alineación de equipos e interfaces GPS; diagrama de proceso}{Alto}{Alto}{Carlos Abarza: revisión de niveles de prueba, controles y evidencias; correspondencia de la figura con el proceso explicado.}
\usoIA{9.3 Alineación con Plan de Trabajo}{Asistente LLM y Codex}{Integración de secuencia de hitos y evidencias}{Alto}{Ninguno}{Carlos Abarza: revisión de secuencia de hitos, ventanas de prueba y condiciones de aceptación y subsanación.}
\usoIA{Formulario T-13}{Asistente LLM y Codex}{Consolidación de gates, evidencia, equipos y perfil nominal/estrés}{Alto}{Ninguno}{Carlos Abarza: revisión de criterios de aprobación, evidencia requerida y consistencia con el plan de calidad.}
\usoIA{Formulario T-17}{Asistente LLM y Codex}{Diseño y depuración de 120 casos; trazabilidad con 42 requisitos; controles de interfaces, carga, recuperación, permisos y adhesión por modalidad}{Alto}{Ninguno}{Carlos Abarza: revisión de claridad de casos, criterios y veredictos; consistencia de referencias con el plan de calidad y la matriz de requerimientos; revisión de modalidades, denominadores, vigencia, revocación y condiciones de capacitación y adhesión.}
